% =========================================================
% IMPLEMENTACIÓN DE LA BASE DE DATOS (tarea O3-03)
% QUÉ VA: cómo el modelo de datos del Obj. 2 (modelo-datos.tex, Anexo 4) quedó
%         en PostgreSQL. Solo lo que sostiene requerimientos clave:
%   1. Migraciones versionadas con Flyway: una por módulo (tabla corta).
%   2. Restricción anticruce de la agenda (RF07): EXCLUDE con btree_gist.
%      Es la evidencia más fuerte de "agenda sin cruces"; mostrar el fragmento
%      SQL (5-6 líneas) como figura o verbatim.
%   3. Separación de datos clínicos e identificación (RNF01, BN3) y
%      anulación en lugar de borrado (RNF04).
%   4. Trazabilidad (RNF02): tablas de historial de Envers (usuario, fecha).
% FUENTE: bellidental-backend/src/main/resources/db/migration y docs/schema.dbml.
% CUIDADO: si el modelo final difiere del diseño, explicarlo en una línea;
%          no reescribir el diagrama ER (ya está en el Anexo 4).
% EXTENSIÓN: 1 tabla + 1 fragmento + 1 párrafo.
% =========================================================

\subsubsection{Base de datos}

El modelo de datos (Figura~\ref{fig:ob2_er_atencion} y Anexo 4) se implementó en PostgreSQL 16 mediante migraciones versionadas con Flyway, una por módulo, de modo que la estructura se reconstruye en cualquier servidor con el mismo orden (Tabla~\ref{tab:ob3_migraciones}).

\begin{table}[H]
\centering
\caption{Migraciones de la base de datos por módulo}
\label{tab:ob3_migraciones}
\fontsize{9}{11}\selectfont
\renewcommand{\arraystretch}{1.2}
\setlength{\tabcolsep}{4pt}
\begin{tabularx}{\linewidth}{|>{\raggedright\arraybackslash}p{3.2cm}|>{\raggedright\arraybackslash}X|>{\raggedright\arraybackslash}p{2.4cm}|}
\hline
\textbf{Módulo} & \textbf{Tablas principales} & \textbf{Requerimiento} \\ \hline
Identidad y acceso & \textit{[usuarios, roles, permisos]} & RNF01 \\ \hline
Pacientes y profesionales & \textit{[pacientes, profesionales, horarios]} & RF06, RF07 \\ \hline
Agenda & \textit{[citas]} & RF07, RF11, RF12 \\ \hline
Clínico & \textit{[historia, antecedentes, atención, diagnósticos, odontograma, tratamientos, recetas]} & RF01 a RF05 \\ \hline
Cobros & \textit{[abonos]} & RF08 \\ \hline
Agente & \textit{[conversaciones, aceptación del aviso, recordatorios]} & RF09 a RF14 \\ \hline
\end{tabularx}
\end{table}

\paragraph{Restricción anticruce}

% Mostrar el fragmento real de la migración de la agenda. Para resaltar
% sintaxis se puede añadir \usepackage{listings} en main.tex; si no, verbatim.
La agenda impide los cruces en la propia base de datos: una restricción de exclusión rechaza dos citas activas del mismo profesional cuyos intervalos se superpongan (RF07). Así, la regla se cumple para la recepción y el agente por igual.

\begin{figure}[H]
\centering
\begin{minipage}{0.9\linewidth}
\small
\begin{verbatim}
-- [Pegar aquí el fragmento EXCLUDE de la migración de la agenda]
EXCLUDE USING gist (
  professional_id WITH =,
  tstzrange(start_at, end_at) WITH &&
) WHERE (status IN (...));
\end{verbatim}
\end{minipage}
\caption{Restricción de exclusión que impide los cruces de citas}
\label{fig:ob3_exclude}
\end{figure}

\textit{[Por completar: 2--3 líneas sobre la separación de datos clínicos e identificación, la anulación en lugar del borrado (RNF04) y las tablas de historial que registran usuario, fecha y hora de cada cambio (RNF02).]}
